%% DVXR Screen — IEEE conference paper (honest write-up).
%% Result tables are auto-generated from the product's evidence layer by
%% scripts/build_paper_tables.py (dvxr.eval.paper.build_product_tables); every number traces to
%% dvxr.serve.evidence.verify_against_scoreboards(), the screeners' subject-held-out aggregation,
%% or a persisted decision curve. Prose states nothing the tables and honesty audit do not support.
\documentclass[conference]{IEEEtran}
\usepackage{booktabs}
\usepackage{graphicx}
\usepackage{amsmath}
\usepackage[hidelinks]{hyperref}

\title{DVXR Screen: A Calibrated, Cross-Subject-Honest, Evidence-Forward\\
System for Multimodal Clinical-Risk Screening from EEG and Wearables}
\author{\IEEEauthorblockN{Abhignya Jagathpally}
\IEEEauthorblockA{DVXR Lab, University of North Texas}}

\begin{document}
\maketitle

\begin{abstract}
Many biosignal classifiers report high accuracy under segment-level splits that leak subject identity,
are never calibrated, and never test whether acting on a prediction is worthwhile. We present DVXR
Screen, a research-grade multimodal clinical-risk screening system that is deliberately conservative on
all three counts. Each task is routed to the representation that wins its own subject-held-out
benchmark: a frozen linear probe over the real LaBraM EEG foundation model for depression screening
from resting EEG, and calibrated band-power physiology for wearable stress. On the Mumtaz resting-EEG
cohort the depression screener reaches a window-level AUROC of 0.961 and a subject-level AUROC of 0.986
($n{=}58$, subject-held-out); we report both granularities because the subject-level number is the
clinically meaningful unit and, here, is not inflated by window pooling. Outputs are Platt-calibrated
with conformal risk intervals, and every screener carries a decision-curve analysis that quantifies the
threshold band over which screening beats the treat-all and treat-none defaults, gated by a bootstrap
lower bound so a chance advantage is not oversold. As a method contribution, a reliability-gated
do-no-harm late fusion beats the proposal's own learned cross-modal fusion on every task and the best
single modality on three of six tasks in the full library; we report where it does not help. We place
every result beside published cross-subject numbers with protocols and DOIs, and refuse comparisons
across mismatched protocols. The system is offline, CPU-only, and deterministic; it is screening and
decision support, not a diagnostic device.
\end{abstract}

\begin{IEEEkeywords}
clinical screening, EEG foundation models, calibration, decision-curve analysis, subject-held-out
cross-validation, multimodal fusion
\end{IEEEkeywords}

\section{Introduction}
Biosignal decoding papers routinely report accuracies above 90\% for tasks such as EEG depression
detection or wearable stress recognition. Three problems recur. First, many of these numbers come from
segment-level train/test splits in which windows from the same subject appear on both sides, so the
reported accuracy overstates performance on a genuinely unseen person. Second, the scores are almost
never calibrated: a model can rank subjects well yet emit probabilities that are systematically wrong,
which matters the moment a threshold is applied. Third, discrimination is reported in isolation from
utility --- an AUROC says nothing about whether acting on the screen is better than the trivial
policies of screening everyone or no one.

DVXR Screen is built to be honest on all three axes at once. Its contributions are: (i) a serving
system that routes each task to the representation that wins its subject-held-out benchmark, headlined
by a frozen probe over the real LaBraM EEG foundation model~\cite{labram}; (ii) reporting at both
window and subject granularity, with the conservative subject-level number placed first; (iii)
Platt calibration and conformal intervals on every output; (iv) decision-curve analysis~\cite{vickers2006dca}
with a bootstrap-gated ``useful'' verdict; and (v) a reliability-gated do-no-harm fusion whose honest
comparison to a learned cross-modal alternative is reported in full, including its failures. Every
number the system surfaces traces to a committed scoreboard, and a blocking honesty audit prevents
excluded results from ever appearing as a claim.

\section{Related Work}
EEG and biosignal foundation models --- LaBraM~\cite{labram}, BENDR~\cite{bendr}, BIOT~\cite{biot},
and the general time-series model MOMENT~\cite{moment} --- provide transferable representations that we
use as frozen feature extractors rather than fine-tuning end-to-end. For depression from EEG, published
cross-subject results are far below the segment-level headlines: MDD-SSTNet reports 65.1\% leave-one-
subject-out (LOSO) accuracy on MODMA~\cite{mddsstnet}, and a GoogleNet CNN drops to 73.3\% under
external validation~\cite{metin2024mdd}, whereas segment-level pipelines such as \cite{yan2022mdd}
report 93.7\% with subject leakage. Wearable stress shows the same pattern: an ensemble trained across
merged corpora reaches roughly 85\% under LOSO~\cite{vos2023stress}, while segment-level WESAD numbers
exceed 94\%~\cite{ghosh2022stress}. For cognitive workload on the PhysioNet mental-arithmetic cohort,
\cite{khanam2023workload} studies the subject-independent setting directly, while \cite{yedukondalu2025workload}
reports 97.4\% on 4\,s segments. Multimodal clinical fusion methods such as MedFuse~\cite{medfuse} and
MedPatch~\cite{medpatch} motivate learned cross-modal fusion; we test such a design and find it loses,
which is why our fusion contribution is a reliability-gated late fusion instead.

\section{Method}
\subsection{Representations and task routing}
For EEG tasks we extract fixed embeddings from a frozen LaBraM~\cite{labram} probe; for wearable and
peripheral-physiology tasks we compute band-power features. Rather than assume a single best
representation, the serving layer routes each task to whichever representation wins its own
subject-held-out benchmark. A logistic head is trained on out-of-fold embeddings.

\subsection{Reliability-gated do-no-harm fusion}
Learned cross-modal fusion is attractive but, in our experiments, harmful (Section~\ref{sec:results}).
We instead use a reliability-gated late fusion in the spirit of oracle model selection~\cite{vanderlaan2007superlearner}:
a modality is admitted to the combined predictor only when it does not degrade held-out performance
relative to the current best. This ``do-no-harm'' gate is conservative by construction; we report the
tasks where it wins and the tasks where the universal guarantee fails at small $N$.

\subsection{Calibration, intervals, and utility}
Head scores are mapped to probabilities by Platt scaling fit on out-of-fold predictions, and we report
the expected calibration error (ECE). A 90\% conformal interval is derived from held-out residuals. For
utility we compute decision-curve analysis~\cite{vickers2006dca}: the net benefit
$\mathrm{NB}(p_t)=\tfrac{\mathrm{TP}}{n}-\tfrac{\mathrm{FP}}{n}\cdot\tfrac{p_t}{1-p_t}$ at decision
threshold $p_t$, computed from the same held-out predictions behind the AUROC, beside the treat-all and
treat-none policies. The ``useful'' band is reported only when the peak net-benefit gain's one-sided
95\% bootstrap lower bound stays positive, so a chance crossing is not sold as utility.

\section{Experiments}
Cohorts are Mumtaz 2016 resting EEG for depression~\cite{mumtaz2016}, WESAD for wearable
stress~\cite{wesad}, the PhysioNet mental-arithmetic cohort for cognitive workload~\cite{eegmat}, and
PhysioNet Non-EEG for peripheral-physiology stress~\cite{noneeg}. All evaluation uses repeated
subject-held-out (subject-disjoint) cross-validation; the head is never fit on a test subject.
Window-level AUROC pools every held-out window; subject-level AUROC aggregates each held-out subject's
windows to a single probability and is reported only for single-class-per-subject tasks. The default
path is offline, CPU-only, and deterministic.

\section{Results}
\label{sec:results}
Table~\ref{tab:headline} reports the headline screening results at both granularities with calibration
error. The depression subject-level AUROC (0.986) exceeds its window-level AUROC (0.961) because
averaging a subject's windows denoises the per-subject signal, so the headline is not an artifact of
window pooling; the interval is nonetheless wide at $n{=}58$.
\IfFileExists{tables/product_headline.tex}{\input{tables/product_headline.tex}}{\textbf{run scripts/build\_paper\_tables.py}}

Table~\ref{tab:fusion} contrasts the winning method for each task with the proposal's own learned
cross-modal fusion. The learned fusion loses on every task; the reliability-gated do-no-harm fusion
beats it on four of six tasks in the full library and the best single modality on three of six. This is
a nuanced positive, not a clean sweep: the universal held-out do-no-harm guarantee does not hold at
$N{\le}60$, where it loses to single-modality ECG on the workload task and slips on the near-chance
affect pair. We report this rather than hide it.
\IfFileExists{tables/fusion_contribution.tex}{\begin{table}[t]
\centering
\caption{Reliability-gated do-no-harm late fusion vs the proposal's own learned cross-modal CACMF fusion. The learned fusion loses on every task; do-no-harm gating beats it on 4 of 6 tasks in the full library (a nuanced positive, not a clean sweep -- see text).}
\label{tab:fusion}
\begin{tabular}{llll}
\toprule
Task & Best method (ours) & Learned CACMF & Delta AUROC \\
\midrule
mumtaz\_depression & LaBraM EEG FM = 0.961 & 0.795 & +0.166 \\
wesad\_stress & band-power + GBM = 0.955 & 0.871 & +0.084 \\
eegmat\_workload & ECG autonomic (task best) = 0.740 & 0.635 & +0.105 \\
stress & peripheral physiology = 0.892 & 0.871 & +0.021 \\
\bottomrule
\end{tabular}
\end{table}
}{}

Table~\ref{tab:sota} places our subject-held-out AUROC beside published results on the same or
comparable cohorts, each labeled with its protocol and DOI. Where a published number is cross-subject
(LOSO or subject-independent) it is the honest bar; where it is segment-level it carries subject
leakage and is shown for context only, never as a head-to-head win. Our contribution is a calibrated,
cross-subject-honest, evidence-forward screener, not a leaderboard accuracy on a leaky split.
\IfFileExists{tables/external_sota.tex}{\input{tables/external_sota.tex}}{}

\subsection{Clinical utility}
Table~\ref{tab:utility} summarizes the decision-curve analysis. The depression screener shows a large,
bootstrap-stable net-benefit advantage over both defaults across a wide threshold band; the weaker
workload screener is useful only over a narrow band, which is the honest reflection of its lower
discrimination.
\IfFileExists{tables/clinical_utility.tex}{\input{tables/clinical_utility.tex}}{}

\section{Limitations}
The EEG cohorts are sampled at 64\,Hz (content below 32\,Hz) against LaBraM's 200\,Hz training regime,
so the EEG results are real but under-resourced on sampling rate and would plausibly improve at native
rate. Cohorts are small ($N{\le}60$), so intervals are wide and the do-no-harm floor diverges from
held-out subjects. Serve-time per-subject recalibration is offered but is an honest negative on these
small cohorts --- on WESAD's eight subjects it worsens ECE (0.156 vs.\ 0.124) --- so it is off by
default. Several results are explicitly excluded from any product claim and blocked by the honesty
audit: affective decoding on DEAP~\cite{deap} (at chance), the learned cross-modal fusion as a win (it
loses on all six tasks), the language-model layer as a predictor (explanation-only), MIMIC
mortality~\cite{mimic}, and an earlier diabetes proxy with label leakage. This is screening and
decision support, never a diagnosis.

\section*{Data Availability}
All cohorts are public research datasets (Mumtaz 2016~\cite{mumtaz2016}, WESAD~\cite{wesad}, PhysioNet
mental-arithmetic~\cite{eegmat}, PhysioNet Non-EEG~\cite{noneeg}); no patient data are used. Result
tables regenerate from the released code via \texttt{scripts/build\_paper\_tables.py}.

\section*{Ethics}
Secondary analysis of de-identified public research cohorts. The system is offline with no PHI egress
and is not a medical device; a positive screen is a prompt to consult a qualified clinician, not a
conclusion.

\section*{Author Contributions}
A.\ Jagathpally: conceptualization, methodology, software, analysis, and writing.

\section*{Conflicts of Interest}
The author declares no competing interests.

\section*{Funding}
No external funding.

\section*{AI Usage Disclosure}
An AI coding assistant was used for software implementation and for drafting and copy-editing this
manuscript. All experimental design, results, and their interpretation are the author's; every reported
number is machine-verified against a committed scoreboard or a persisted model artifact.

\bibliographystyle{IEEEtran}
\bibliography{references}
\end{document}
